\section{Introduction}

\subsection{Motivation}

\HE{} is organized around the following security shape:
\[
\begin{gathered}
\text{public handles}\\
\text{and verifier-facing artifacts}
\end{gathered}
\quad\not\Rightarrow\quad
\text{the operative Hodge--Laplacian kernel}.
\]
The private recovery object is the harmonic cochain selected by a hidden Hodge
operator, instead of a scalar or a vector secret in a published lattice or algebraic object. 
Public verification therefore needs finite artifacts, while security
must not collapse to recovering the private representative or the hidden
operator.

\subsection{Contributions}

This paper develops the following contributions.
\begin{enumerate}[label=(\roman*)]
  \item The Harmonic Cochain Problem (HCP): recovery of the harmonic cochain
    selected by a hidden Hodge operator.
  \item A Hodge-structured noise model with cochain decomposition
    \(b = h + e\), where \(h \in \ker(\Delta)\) and
    \(e \in \image(\delta) \oplus \image(\delta^\ast)\).
  \item A CVP comparison in which an explicit lattice is realized as a kernel
    \(\ker(\Delta)\), a harmonic solver yields a lattice vector, and the
    closest-harmonic contract recovers the closest-vector predicate.
  \item A hidden-operator ambiguity theorem: publicly equivalent Hodge operators
    may disagree on harmonicity, so a correct explicit-lattice CVP oracle does
    not uniformly solve public-data HCP.
  \item A Lean formalization of the definitions, reductions, projection chain,
    gates, and theorem boundaries.
  \item A projection chain through Plucker shadows and cochain commitments to a
    canonical Reed--Solomon endpoint with an MCA bad-challenge bound.
\end{enumerate}
